% chapters/requirements/rf_users.tex
% RF-2: Gestión de usuarios y organizaciones
% ============================================================================

\item \textbf{Gestión de usuarios y organizaciones.} \label{rf:users}
El sistema se estructura en torno a organizaciones como contenedor principal
\dfn{multi-tenant}. Una organización representa una institución (universidad, instituto)
en modo instalación única, o un cliente en modo SaaS. Los usuarios pertenecen a una única
organización; un mismo individuo en otra organización se considera un usuario distinto
(otra cuenta). Los usuarios no se eliminan físicamente: la baja se modela mediante
desactivación lógica (\texttt{is\_active=False}) para preservar la integridad referencial
con calificaciones, anotaciones y exámenes históricos. Los gestores
(\texttt{is\_staff=True}) son los únicos que pueden crear, modificar y desactivar usuarios
dentro de su organización. Los roles académicos se definen exclusivamente en el contexto
de una asignatura mediante \texttt{SubjectMembership} (véase \cref{rf:subjects}).

\begin{functional}

    % RF-2.1 ----------------------------------------------------------------
    \item \textbf{Creación de organización.} \label{rf:users:create_org}
    El sistema debe permitir la creación de nuevas organizaciones.

    \textbf{Entrada:} petición POST al recurso de organizaciones, accesible solo por el
    superadministrador del sistema, con los datos: nombre de la organización, subdominio
    único, plan contratado (FREE o PREMIUM) y datos del gestor inicial.

    \textbf{Proceso:} el sistema valida la unicidad del subdominio, crea la organización,
    crea el usuario gestor inicial (\texttt{is\_staff=True}, \texttt{is\_active=True}) y le
    envía un correo de bienvenida con su contraseña generada. Se registra la operación en
    el \textit{log} de auditoría global.

    \textbf{Salida:} respuesta con los datos de la organización creada y código
    \acs{http}~201.

    % RF-2.2 ----------------------------------------------------------------
    \item \textbf{Creación individual de usuario.} \label{rf:users:create}
    El sistema debe permitir a un gestor crear un usuario dentro de su organización.

    \textbf{Entrada:} petición POST al recurso de usuarios con los datos del nuevo usuario:
    nombre, apellidos, \acs{dni}, identificador en la organización (\acs{nia}), correo
    electrónico y otros datos opcionales.

    \textbf{Proceso:} el sistema valida la unicidad del correo electrónico, el \acs{nia} y
    el \acs{dni} dentro de la organización. Cifra el \acs{dni} con AES-256-GCM antes de
    almacenarlo (véase \cref{rf:sec:dni}). Genera una contraseña segura aleatoria, la
    almacena como \textit{hash} con \textit{salt} y encola un correo electrónico de
    bienvenida con la contraseña en texto plano. El usuario se crea con
    \texttt{is\_active=True} y \texttt{is\_staff=False}. Se registra la operación en el
    \textit{log} de auditoría.

    \textbf{Salida:} respuesta con los datos del usuario creado (sin contraseña ni
    \acs{dni} descifrado) y código \acs{http}~201.

    % RF-2.3 ----------------------------------------------------------------
    \item \textbf{Importación masiva de usuarios.} \label{rf:users:import}
    El sistema debe permitir a un gestor crear o reactivar múltiples usuarios
    simultáneamente a partir de un fichero estructurado.

    \textbf{Entrada:} petición POST con un fichero en formato \acs{csv} o \acs{json} que
    contenga los datos de cada usuario (nombre, apellidos, \acs{dni}, \acs{nia}, correo,
    etc.).

    \textbf{Proceso:} el sistema parsea el fichero, valida cada registro individualmente y
    procesa la importación de forma transaccional. Para cada registro: si el usuario no
    existe en la organización, se crea con \texttt{is\_active=True}; si ya existe pero
    está inactivo, se reactiva y se actualizan sus datos personales; si ya existe y está
    activo, se actualizan sus datos. Para cada usuario se genera una contraseña segura y se
    encola el correo de bienvenida. Se registra la operación en el \textit{log} de
    auditoría con el número de usuarios creados, reactivados y actualizados.

    \textbf{Salida:} respuesta con un informe del resultado: número de usuarios creados,
    reactivados, actualizados y omitidos por error, con detalle de los errores de
    validación por fila.

    % RF-2.4 ----------------------------------------------------------------
    \item \textbf{Modificación de datos de usuario.} \label{rf:users:update}
    El sistema debe permitir a un gestor modificar los datos de cualquier usuario de su
    organización.

    \textbf{Entrada:} petición PATCH al recurso del usuario con los campos a modificar
    (nombre, apellidos, correo).

    \textbf{Proceso:} el sistema valida los datos actualizados, mantiene la unicidad de los
    campos restringidos dentro de la organización y aplica los cambios. Se registra el
    cambio en el \textit{log} de auditoría con los valores anteriores y nuevos.

    \textbf{Salida:} respuesta con los datos actualizados del usuario.

    % RF-2.5 ----------------------------------------------------------------
    \item \textbf{Desactivación de usuario.} \label{rf:users:deactivate}
    El sistema debe permitir a un gestor dar de baja a un usuario sin eliminar su registro
    de la base de datos.

    \textbf{Entrada:} petición PATCH al recurso del usuario estableciendo
    \texttt{is\_active=False}.

    \textbf{Proceso:} el sistema marca al usuario como inactivo. Un usuario inactivo no
    puede autenticarse ni acceder a ningún recurso, pero sus datos y relaciones históricas
    (calificaciones, membresías en cursos archivados) se mantienen intactos. Se registra la
    operación en el \textit{log} de auditoría.

    \textbf{Salida:} respuesta de confirmación con los datos del usuario actualizado.

    % RF-2.6 ----------------------------------------------------------------
    \item \textbf{Gestión de gestores.} \label{rf:users:staff}
    El sistema debe permitir a un gestor existente designar o revocar el rol de gestor a
    otros usuarios de la organización.

    \textbf{Entrada:} petición PATCH al recurso del usuario modificando el campo
    \texttt{is\_staff}.

    \textbf{Proceso:} el sistema verifica que el usuario que realiza la operación es gestor
    y pertenece a la misma organización. El usuario designado como gestor obtiene permisos
    de administración sobre la organización. Los gestores no se ven afectados por la
    desactivación automática en los cambios de curso (véase \cref{rf:courses:transition}).
    Se registra la operación en el \textit{log} de auditoría.

    \textbf{Salida:} respuesta con los datos actualizados del usuario.

    % RF-2.7 ----------------------------------------------------------------
    \item \textbf{Generación y envío de contraseña segura.} \label{rf:users:password_gen}
    El sistema debe generar automáticamente una contraseña segura al crear un usuario y
    enviarla por correo electrónico.

    \textbf{Entrada:} evento interno disparado tras la creación de un usuario.

    \textbf{Proceso:} el sistema genera una contraseña aleatoria que cumpla requisitos
    mínimos de complejidad (longitud, mayúsculas, minúsculas, dígitos y caracteres
    especiales). La contraseña se almacena como \textit{hash} con \textit{salt} mediante
    Argon2 o bcrypt. Se encola una tarea asíncrona en la cola de notificaciones que envía
    un correo al usuario con la contraseña en texto plano y las instrucciones de primer
    acceso.

    \textbf{Salida:} correo electrónico enviado al usuario. La contraseña en texto plano no
    se almacena en el sistema tras su envío.

    % RF-2.8 ----------------------------------------------------------------
    \item \textbf{Restablecimiento de contraseña por gestor.} \label{rf:users:password_reset}
    El sistema debe permitir a un gestor forzar el restablecimiento de la contraseña de un
    usuario de su organización.

    \textbf{Entrada:} petición POST al recurso de restablecimiento de contraseña indicando
    el identificador del usuario.

    \textbf{Proceso:} el sistema verifica que el usuario pertenece a la misma organización.
    Genera una nueva contraseña segura, actualiza el \textit{hash} almacenado, invalida
    todos los \textit{tokens} activos del usuario (añadiendo los \textit{refresh tokens} a
    la lista negra en Redis) y encola el envío de la nueva contraseña por correo. Se
    registra la operación en el \textit{log} de auditoría.

    \textbf{Salida:} respuesta de confirmación con código \acs{http}~200.

    % RF-2.9 ----------------------------------------------------------------
    \item \textbf{Modificación del correo electrónico propio.} \label{rf:users:own_email}
    El sistema debe permitir a cualquier usuario autenticado modificar su propia dirección
    de correo electrónico.

    \textbf{Entrada:} petición PATCH al recurso del perfil propio con el nuevo correo
    electrónico.

    \textbf{Proceso:} el sistema valida que el nuevo correo no esté ya registrado en la
    misma organización y aplica el cambio. Se registra la modificación en el \textit{log}
    de auditoría.

    \textbf{Salida:} respuesta con los datos actualizados del perfil.

    % RF-2.10 ---------------------------------------------------------------
    \item \textbf{Consulta del perfil propio.} \label{rf:users:own_profile}
    El sistema debe permitir a cualquier usuario autenticado consultar sus datos personales,
    las asignaturas en las que participa y su rol en cada una.

    \textbf{Entrada:} petición GET al recurso del perfil propio.

    \textbf{Proceso:} el sistema recupera los datos del usuario autenticado, incluyendo
    nombre, apellidos, correo, preferencias de notificación y la lista de asignaturas del
    curso activo con el rol que ostenta en cada una (obtenido de
    \texttt{SubjectMembership}). El \acs{dni} se incluye descifrado solo si el usuario
    consulta su propio perfil o es gestor.

    \textbf{Salida:} respuesta con los datos del perfil.

    % RF-2.11 ---------------------------------------------------------------
    \item \textbf{Gestión de preferencias de notificación.} \label{rf:users:notif_prefs}
    El sistema debe permitir a cualquier usuario activar o desactivar la recepción de
    notificaciones por correo electrónico.

    \textbf{Entrada:} petición PATCH al recurso del perfil propio con el campo de
    preferencia de notificaciones por correo.

    \textbf{Proceso:} el sistema actualiza la preferencia. Cuando las notificaciones por
    correo están desactivadas, las notificaciones solo se almacenan en la base de datos
    (notificaciones \textit{in-app}) sin enviar correo espejo.

    \textbf{Salida:} respuesta con la preferencia actualizada.

    % RF-2.12 ---------------------------------------------------------------
    \item \textbf{Listado y búsqueda de usuarios.} \label{rf:users:list}
    El sistema debe permitir a un gestor listar y buscar usuarios de su organización con
    filtros avanzados.

    \textbf{Entrada:} petición GET al recurso de usuarios con parámetros opcionales de
    filtrado (nombre, apellidos, \acs{nia}, correo, estado activo o inactivo) y paginación.

    \textbf{Proceso:} el sistema aplica los filtros indicados sobre el conjunto de usuarios
    de la organización y devuelve los resultados paginados. Se permite la búsqueda parcial
    por texto en nombre y apellidos.

    \textbf{Salida:} respuesta con la lista paginada de usuarios que cumplen los criterios,
    incluyendo los metadatos de paginación (total, página actual, tamaño de página).

    % RF-2.13 ---------------------------------------------------------------
    \item \textbf{Exportación de usuarios.} \label{rf:users:export}
    El sistema debe permitir a un gestor exportar la lista de usuarios de su organización
    en formato estructurado.

    \textbf{Entrada:} petición GET al recurso de exportación de usuarios con parámetros
    opcionales de formato (\acs{csv} o \acs{json}) y filtros.

    \textbf{Proceso:} el sistema genera un fichero con los datos de los usuarios que cumplen
    los filtros indicados. Los datos sensibles (\acs{dni}) se incluyen descifrados en la
    exportación. Se registra la operación en el \textit{log} de auditoría.

    \textbf{Salida:} respuesta con el fichero generado como descarga.

\end{functional}
